\documentclass[10pt,a4paper]{amsart}
\usepackage[margin=19mm]{geometry}
\usepackage{amsmath,amssymb,mathtools}
\usepackage[scr=rsfs,cal=euler]{mathalfa}
\usepackage{xfrac}
\usepackage[unicode,hidelinks,bookmarks=false]{hyperref}
\newcommand{\ie}{i.e.\ }
\newcommand{\cf}{cf.\ }
\newcommand{\resp}{respectively\ }
\newcommand{\Subsection}{\subsection}
\providecommand{\nobreakdash}{\nobreak\hbox{-}}
\setlength{\emergencystretch}{2em}
\multlinegap0pt
\usepackage{amssymb}
\usepackage{xcolor}
\usepackage{bm}
\usepackage{mathtools}
\usepackage{tikz}
\usepackage{dsfont}
\usepackage[shortlabels]{enumitem}
\newtheorem{lem}{Lemma}
\title{Eigenfunction equivalence for the fractional Laplace-Beltrami operator and the classical Helmholtz equation}
\author{Saumyajit Das, Susovan Pramanik}
\date{Source: arXiv:2608.09211v1 (CC BY 4.0)}
\begin{document}
\maketitle

\noindent\emph{Source and licence.} Eigenfunction equivalence for the fractional Laplace-Beltrami operator and the classical Helmholtz equation. Version arXiv:2608.09211v1 (CC BY 4.0), distributed by arXiv under the Creative Commons Attribution 4.0 International licence, http://creativecommons.org/licenses/by/4.0/, as recorded in the arXiv licence metadata for this exact version and shown on the abstract page https://arxiv.org/abs/2608.09211v1. Full licence notice: \texttt{licenses/arxiv-2608.09211-CC-BY-4.0.txt}.\par\smallskip

\noindent\textit{Editorial scope.} Counted unit: the article's lem calculus (source0.tex lines 590--609). The article's complete original proof of it is delivered with it (source0.tex lines 610--650, one \texttt{proof} environment of 41 lines). No mathematics is written, repaired or re-derived: the statement and the proof are the article's own lines, in the article's own notation, and no retained line is substituted or re-flowed.  Retained uncounted as the source-local context the proof needs: lines 549--551.  Nothing was omitted from any retained span, so the package carries no omission record.  No retained line contains a citation, so the wrapper carries no bibliography and no bibliography entry.  No retained line contains a figure environment, an image-inclusion command or drawing code, so no image is delivered and the package declares no asset dependency.  Editorial formatting only, in addition: an AMS \texttt{amsart} preamble that restates the article's own declarations and macros used by the retained lines, namely the environments \texttt{lem} on separate counters, together with the standard packages those lines need.  The retained environments print in this document as Lemma 1 in order of appearance, whereas the article prints the same items with the numbering of its own full document.  The complete native source of the article as distributed in the arXiv e-print for this exact version is bundled unchanged as \texttt{sources/207224/source0.tex}.
	\begin{lem}\label{decay of derivatives}
		Let $l\in\mathbb{N}\cup\{0\}$, then $\displaystyle{\left|D^l(b_k)(\lambda)\right|\lesssim |\lambda|^{1-s-l}}$, as $\lambda\to\infty$.
	\end{lem}

	\begin{lem}\label{towards functional calculus}
		Let $\bigl(g^{ij}(x)\bigr)_{i,j=1}^{d}$ be a uniformly elliptic, smooth, and bounded matrix-valued function whose derivatives of all orders are bounded. Then, for every pair of multi-indices $\alpha,\beta\in (\mathbb{N}\cup{0})^{d}$, the estimate
		\[
		\left|\partial_x^{\alpha}\partial_{\xi}^{\beta} b_k\left(\sum_{i,j=1}^d g^{ij}(x)\xi_i\xi_j\right)\right|
		\lesssim
		\langle \xi\rangle^{2-2s-|\beta|}
		\]
		holds uniformly for all $x,\xi\in\mathbb{R}^{d}$. Here,
		\[
		\langle \xi\rangle
		:=
		\left(1+\sum_{i=1}^{d}|\xi_i|^2\right)^{1/2}
		\]
		denotes the Japanese bracket. Moreover, for a multi-index
		$r=(r_1,\ldots,r_d)\in(\mathbb{N}\cup{0})^d$,
		we define
		\[
		|r|:=\sum_{i=1}^{d} r_i.
		\]
	\end{lem}

	\begin{proof}
		Thanks to uniform ellipticity of the matrix valued function $g^{ij}(x)$, we have that
		\[
		\theta |\xi|^2 \leq \sum_{i,j=1}^d g^{ij}(x)\xi_i\xi_j\leq \Theta |\xi|^2,
		\qquad x,\xi\in\mathbb{R}^d,
		\]
		where $\theta, \Theta$ are some positive constants. Let $q(x,\xi):=\sum_{i,j=1}^d g^{ij}\xi_i\xi_j$. We use Fa\'a di Bruno's formula to compute the mixed derivatives. We have that  
		\[
		\partial^{\beta}_{\xi} b_k(q)= \sum_{m=1}^{|\beta|} (D^m b_k)(q) \prod_{i=1}^m \partial_{\xi}^{\mu_i} q, 
		\]
		where $\mu_i\in (\mathbb{N}\cup\{0\})^d$ for all $i=1,\cdots,m$ and $\displaystyle{\sum_{i=1}^m |\mu_i|=|\beta|}$. Since $q(x,\xi)$ is a polynomial of order two, $\partial_{\xi}^{\mu_i}q$ vanishes where $|\mu_i|\geq 3$. Hence $\displaystyle{\prod_{i=1}^m \partial_{\xi}^{\mu_i} q}$ is a polynomial of degree $\displaystyle{\sum_{i=1}^m (2-|\mu_i|)}=2m-|\beta|$ with coefficients that are smooth bounded functions of $x$. Moreover, all derivatives of these coefficients are uniformly bounded. Next we use Leibnitz rule. 
		\[
		\partial_{x}^{\alpha}\partial_{\xi}^{\beta} b_k(q)= \sum_{m=1}^{|\beta|} \left( \sum_{\nu\leq \alpha} \begin{pmatrix}
			\alpha\\ \nu
		\end{pmatrix} \partial^{\nu}_{x}(D^m b_k)(q) \partial^{\alpha-\nu}_{x}\left(\prod_{i=1}^m \partial_{\xi}^{\mu_i} q\right)\right),
		\]
		where $\nu\in(\mathbb{N}\cup\{0\})^d$, $\nu\leq \alpha$ means $\nu\leq \alpha_i$ for all $i=1,\cdots, d$ and $\displaystyle{\begin{pmatrix}
				\alpha\\ \nu
			\end{pmatrix}=\prod_{i=1}^d \begin{pmatrix}
				\alpha_i\\ \nu_i
		\end{pmatrix}}$. We use Fa\`a di Bruno's formula once again. It yields
		\begin{align}\label{intermediate 1}
			\partial_{x}^{\alpha}\partial_{\xi}^{\beta} b_k(q)= \sum_{m=1}^{|\beta|} \left( \sum_{\nu\leq \alpha} \begin{pmatrix}
				\alpha\\ \nu
			\end{pmatrix} \left(\sum_{j=1}^{|\nu|} (D^{m+j} b_k)(q) \prod_{l=1}^{j} \partial_x^{\delta_l} q(x)\right) \partial^{\alpha-\nu}_{x}\left(\prod_{i=1}^m \partial_{\xi}^{\mu_i} q\right)\right),
		\end{align}
		where $\delta_l\in (\mathbb{N}\cup\{0\})^d$ for all $l=1,\cdots,j$ and $\displaystyle{\sum_{l=1}^j |\delta_l|=|\nu|}$. Note that the function $q$ is positive thanks to uniform ellipticity. We can further say that the term $\displaystyle{\partial^{\alpha-\nu}_{x}\left(\prod_{i=1}^m \partial_{\xi}^{\mu_i} q\right)}$ is a polynomial (in $\xi$) of degree $2m-|\beta|$ with smooth and bounded coefficients in `$x$' and similarly the term $\displaystyle{\prod_{l=1}^{j} \partial_x^{\delta_l} q(x)}$ is a polynomial (in $\xi$) of degree $\displaystyle{\sum_{l=1}^j 2=2j}$ with smooth bounded coefficients in `$x$'. Moreover, thanks to Lemma \ref{decay of derivatives}, we have that 
		\[
		\left| (D^{m+j} b_k)(q)\right| \lesssim | q|^{1-s-m-j}, \quad \mbox{as} \ q\to \infty. 
		\]
		Thanks to uniform ellipticity again we can write the above expression as 
		\[
		\left| (D^{m+j} b_k)(q)\right| \lesssim | \xi|^{2-2s-2m-2j}, \quad \mbox{as} \ \xi\to \infty.
		\]
		Hence, from \eqref{intermediate 1}, we conclude that
		\[
		\left|\partial_x^{\alpha}\partial_{\xi}^{\beta} b_k\left(\sum_{i,j=1}^d g^{ij}(x)\xi_i\xi_j\right)\right|
		\lesssim
		\langle \xi\rangle^{2-2s-2m-2j+2j+2m-|\beta|}=  \langle \xi\rangle^{2-2s-|\beta|}, \quad \forall\, x, \xi\in\mathbb{R}^d.
		\]
	\end{proof}

\end{document}
